\documentclass{article}

\usepackage[utf8]{inputenc}
\usepackage[T1]{fontenc}
\usepackage{amsmath,amssymb,amsthm}
\usepackage{mathtools}
\usepackage{hyperref}
\usepackage{booktabs}

\makeatletter
\newcommand{\citep}[2][]{%
  \begingroup
  \def\@tmpkey{#2}%
  \ifx\@tmpkey\@empty\else
    (\ifx\@empty#1\@empty\else #1, \fi
     \@nameuse{bibkey@\@tmpkey})%
  \fi
  \endgroup
}
\expandafter\def\csname bibkey@lasser2026exo\endcsname{Exogenous~Verification}
\expandafter\def\csname bibkey@lasser2026phase\endcsname{Phase~Redundancy}
\expandafter\def\csname bibkey@lasser2026micro\endcsname{Microfoundation}
\makeatother

\newtheorem{definition}{Definition}
\newtheorem{lemma}{Lemma}
\newtheorem{theorem}{Theorem}
\newtheorem{remark}{Remark}
\newtheorem{assumption}{Assumption}

\newcommand{\Aadv}{A_{\mathrm{adv}}}
\newcommand{\deltaadv}{\delta_{\mathrm{adv}}}
\newcommand{\deltastar}{\delta_{*}}
\newcommand{\rhogap}{\rho_{\mathrm{gap}}}
\newcommand{\Pproxy}{P}
\newcommand{\Ttruth}{T}
\newcommand{\epsgap}{\varepsilon_{\mathrm{gap}}}
\newcommand{\epsnonres}{\varepsilon_{\mathrm{gap}}^{\mathrm{nonres}}}
\newcommand{\epsfloor}{\varepsilon_{\mathrm{floor}}^{\mathrm{res}}}

\newcommand{\Tbeta}{T_{\beta}}

\title{Lemma 6 ($h_{\mathrm{detect}}$ intensivity): Full Proof Draft v3}
\author{Paper 10 working draft for codex re-review}
\date{}

\begin{document}

\maketitle

\section{Goal and changes from v2}

V2 fixed v1's two structural problems
(integrated-increment $h_{\mathrm{detect}}$; missing monitoring
clock) but Round~B identified one significant problem and two
issues:

\textbf{Round B significant problem.}  V2's Step~6 used
$T_{\mathrm{cascade}}$ as an upper bound on $T_{\mathrm{detect}}$,
but \citep[Phase~Redundancy]{lasser2026phase}'s metastable-lifetime
result gives only the \emph{lower} bound
$T_{\mathrm{cascade}} \geq \tau_{\mathrm{meta}}$ — useful for
lead-time-before-cascade but not as the integration horizon for
$\rhogap \cdot T$.

\textbf{Round B issues.}  Q1: discrete-time boundary condition
$\epsgap(t_0) \leq h_{\mathrm{static}}$ was not rigorous (one-step
jump was ``absorbed'' without explicit convention).  Q2: fixed-step
event-rate normalization needed an explicit named clause, not
implicit in (C5).

\textbf{Changes from v2 (Round B fixes):}

\begin{itemize}
\item \textbf{Significant problem fix
  ($T_{\mathrm{cascade}} \to \Tbeta$).}  Introduced
  $\Tbeta$ as the SPRT high-probability detection \emph{quantile}:
  $\Pr[T_{\mathrm{detect}} \leq \Tbeta] \geq 1-\beta$, where
  $\Tbeta := (\log(1/\alpha) + \log(1/\beta)) / (\kappa \deltastar)$
  follows directly from Lemma~4's Wald--Hoeffding tail with
  KL floor $\deltastar$.  This is a legitimate upper bound on
  $T_{\mathrm{detect}}$, intensive in $|P|$ since both $\kappa$ and
  $\deltastar$ are intensive.  The integrated bound becomes
  $h_{\mathrm{detect}} \leq h_{\mathrm{static}} + \rhogap \cdot
  \Tbeta$ with probability at least $1-\beta$.
  Lead-time-before-cascade is then a separate composition with
  Lemma~5d: $\Tbeta \leq \tau_{\mathrm{meta}} \leq T_{\mathrm{cascade}}$
  with probability at least $1-\beta'$.
\item \textbf{Q1 fix (discrete boundary).}  Explicit convention:
  $t_0$ is defined as the last SPRT exposure step before the
  Layer~1$\to$Layer~2 crossing (i.e., $t_0 := t_0^-$ on the SPRT
  exposure clock), and $T_{\mathrm{detect}}$ is counted from $t_0$
  inclusive of the crossing step.  This makes
  $\epsgap(t_0) \leq h_{\mathrm{static}}$ rigorous without an
  ``absorbed'' single-step buffer.
\item \textbf{Q2 fix (clock comparability).}  Added (C11.CLK)
  named sub-clause: detection and cascade times are expressed in
  the SPRT exposure clock, with a calibrated high-probability lower
  bound on event count before cascade.  This eliminates the
  implicit ``average inter-event interval'' assumption and handles
  bursty event-rate deployments.
\item \textbf{Q4 reinforcement.}  Explicit statement that $\rhogap$
  is calibrated in the $\epsgap$ metric (the proxy-truth gap, not
  the alignment-property error) and excludes $K_{\mathrm{Lip}}$.
\end{itemize}

\textbf{Carried forward from v2 (unchanged in v3):}
\begin{itemize}
\item Supremum-form $h_{\mathrm{detect}}(\theta) := \sup_t
  \epsgap(t)$ matching Theorem~1's Layer~2 application.
\item Renaming $\rho_g \to \rhogap$ to clarify gap-growth-rate
  semantics.
\item Union-class KL floor $\deltastar = \min(\deltaadv,
  \delta_{\mathrm{adv}}^{\mathrm{env}})$ for $\Aadv \cup
  \Aadv^{\mathrm{env}}$.
\item (C11) explicitly framed for promotion to theorem-level
  condition.
\item Source-citation correction:
  \citep[Microfoundation]{lasser2026micro}'s ``$T$ failure modes''
  remark does not establish the per-step gap-growth rate; (C11) is
  a new Paper~10 operational assumption.
\end{itemize}

\section{Setup: the monitoring clock and $h_{\mathrm{detect}}$}

Theorem~\ref{thm:deployment_safety}'s Layer~2 covers adversarial
strategies $q \in \Aadv \cup \Aadv^{\mathrm{env}}$ during the
detection window
$[t_0, t_0 + T_{\mathrm{detect}}]$.

\paragraph{Boundary convention.}  We define $t_0$ as the
\emph{last SPRT exposure step} at which the deployment is still
within Layer~1's safe region; equivalently, $t_0$ is the SPRT step
immediately preceding the first crossing event ($t_0 := t_0^-$ in
the v2 notation).  $T_{\mathrm{detect}}$ is the number of SPRT
steps from $t_0$ until the SPRT log-likelihood ratio crosses the
detection threshold (Lemma~\ref{lem:4}); the crossing step itself
is counted within $T_{\mathrm{detect}}$.

Under this convention, at $s = s_0 = t_0$ the Layer~1 bound applies
exactly: $\epsgap(t_0) \leq h_{\mathrm{static}}(\theta)$.  The
single-step jump that first crosses Layer~1 occurs at $s_1$, and is
bounded by $\rhogap(\deltastar)$ via (C11).  No ``absorbed buffer''
is required.

\paragraph{Detection-window bound.}
Layer~2 applies $h_{\mathrm{detect}}$ as a point-in-time bound:
\[
  |g(\Ttruth) - g(\Pproxy)|(s) \;\leq\; \mathrm{Lip}(g) \cdot
  h_{\mathrm{detect}}(\theta)
  \quad \text{for all } s \in [t_0, t_0 + T_{\mathrm{detect}}].
\]
Therefore $h_{\mathrm{detect}}$ must bound $\epsgap(s)$ at every
$s$ in the window:
\[
  h_{\mathrm{detect}}(\theta) \;:=\; \sup_{s \in
  [t_0, t_0 + T_{\mathrm{detect}}]} \epsgap(s).
\]

\paragraph{The monitoring clock and (C11.CLK) clock comparability.}
The SPRT machinery (Paper~5) operates on a per-event clock: each
ledger entry triggers an SPRT update, and the per-step KL divergence
drives detection.  All times in this lemma — $t_0$,
$T_{\mathrm{detect}}$, $\Tbeta$, and the cascade-time references
$\tau_{\mathrm{meta}}$, $T_{\mathrm{cascade}}$ — are expressed in
this SPRT exposure clock.

\begin{assumption}[(C11.CLK): Clock comparability]
\label{ass:C11CLK}
The deployment's event-counting process satisfies a calibrated
high-probability lower bound on the number of SPRT exposure steps
elapsing before any cascade event: there exist deployment-class
constants $N_{\min}, q_{\min} > 0$ such that
$\Pr[N_{\mathrm{events}}(\tau_{\mathrm{meta}}) \geq N_{\min}] \geq
q_{\min}$, where $N_{\mathrm{events}}(\cdot)$ is the event count up
to wall-clock time $\cdot$.  Equivalently, the SPRT exposure clock
is calibrated to wall-clock cascade time with a known floor on
event throughput.
\end{assumption}

\textbf{Why (C11.CLK) is needed.}  Without it, bursty event-rate
deployments could in principle produce $T_{\mathrm{detect}}$ in
SPRT steps yet have wall-clock detection time exceeding cascade
time.  (C11.CLK) eliminates this gap by making the event-counting
process itself calibrated.

\section{Bounded gap-growth-rate condition (C11)}

\begin{assumption}[(C11): Bounded gap-growth rate per SPRT step]
\label{ass:C11}
For any adversarial strategy $q \in \Aadv \cup
\Aadv^{\mathrm{env}}$ producing channel deviations bounded by
$\deltastar = \min(\deltaadv, \delta_{\mathrm{adv}}^{\mathrm{env}})$
(the Lemmas~5b and~5e KL floors), the per-SPRT-step proxy-truth
gap growth is bounded:
\[
  (\Delta \epsgap)_n \;:=\; \epsgap(s_n) - \epsgap(s_{n-1}) \;\leq\;
  \rhogap(\deltastar),
\]
where $\rhogap(\deltastar)$ is a deployment-class constant
depending on $\mathcal{D}$'s operational parameters (channel
sensitivity-to-gap mapping, event-rate normalization) but
independent of $|P|$.

\textbf{Continuous-time variant.}  Equivalently, in continuous
SPRT exposure clock $s$:
\[
  \frac{d \epsgap}{ds} \;\leq\; \rhogap(\deltastar).
\]
The discrete and continuous forms agree under fixed step duration
or bounded event-rate normalization.
\end{assumption}

\textbf{Operational interpretation.}  Channel deviations bounded by
$\deltastar$ (else SPRT fires faster) limit how much $\epsgap$ can
grow per SPRT step.  The constant $\rhogap$ characterizes the
deployment's channel-deviation-to-gap-growth coupling.

\textbf{$\rhogap$ is calibrated in the $\epsgap$ metric only.}
$\rhogap$ is a bound on the per-step proxy-truth gap process
$\epsgap = $ (Microfoundation's relative sup-norm), \emph{before}
applying any alignment-property mapping $g$.  In particular,
$\rhogap$ excludes the Lipschitz constant $K_{\mathrm{Lip}}$ of $g$;
the Layer~2 operational form multiplies $h_{\mathrm{detect}}$ by
$K_{\mathrm{Lip}}$ separately via (C6).  This eliminates any
double-counting between (C11) gap-dynamics and (C6) alignment-side
smoothness.

\textbf{(C11) is new content.}  \citep[Microfoundation]{lasser2026micro}'s
``$T$ failure modes'' remark notes adequacy/failure modes of the
operational truth $T$ but does not establish a formal per-step
gap-growth-rate bound.  (C11) is a Paper~10 operational assumption,
analogous to (C7) bounded co-evolution and SA1 HHI surrogate
adequacy: deployment-class condition with empirical testability,
not a derived consequence of source-paper machinery.

\section{Statement of Lemma 6 (v3)}

\begin{lemma}[$h_{\mathrm{detect}}$ intensivity]
\label{lem:6}
Under Theorem~\ref{thm:deployment_safety}'s conditions (C5)
continuous SPRT monitoring, (C11) bounded gap-growth rate, and
(C11.CLK) clock comparability, combined with Lemma~\ref{lem:4}
(SPRT Wald--Hoeffding tail bound on $T_{\mathrm{detect}}$) and the
Layer~1 intensivity result Lemma~\ref{lem:1}:

\textbf{Detection-window supremum bound.}  Define the SPRT
high-probability detection quantile
\[
  \Tbeta \;:=\;
  \frac{\log(1/\alpha) + \log(1/\beta)}{\kappa \cdot \deltastar},
\]
where $\alpha$ is the Type-I error rate, $\beta$ is the Type-II
error rate (target $\Pr[T_{\mathrm{detect}} > \Tbeta] \leq \beta$),
$\kappa$ is the Hoeffding sub-Gaussian constant, and $\deltastar$
is the union-class KL floor.  Then:

(i) [Tail bound on $T_{\mathrm{detect}}$.]  By Lemma~\ref{lem:4}'s
Wald--Hoeffding bound,
$\Pr[T_{\mathrm{detect}} \leq \Tbeta] \geq 1 - \beta$.

(ii) [Detection-window supremum.]  On the event
$\{T_{\mathrm{detect}} \leq \Tbeta\}$ (probability at least
$1-\beta$):
\[
  h_{\mathrm{detect}}(\theta) \;:=\;
  \sup_{s \in [t_0, t_0 + T_{\mathrm{detect}}]} \epsgap(s)
  \;\leq\; h_{\mathrm{static}}(\theta) + \rhogap(\deltastar) \cdot
  \Tbeta.
\]

(iii) [Lead-time-before-cascade.]  Separately, by
Lemma~\ref{lem:5d}'s composition with
\citep[Phase~Redundancy]{lasser2026phase}'s metastable lifetime,
$\Pr[\Tbeta \leq \tau_{\mathrm{meta}}] \geq 1 - \beta'$ where
$\beta' = \exp(-\kappa \tau_{\mathrm{meta}} \deltastar)$, and
$\tau_{\mathrm{meta}} \leq T_{\mathrm{cascade}}$.

\textbf{Intensivity.}  $h_{\mathrm{detect}}(\theta)$ is intensive
in $|P|$ over the deployment class $\mathcal{D}$:
$h_{\mathrm{static}}(\theta)$ is intensive by Lemma~\ref{lem:1};
$\rhogap(\deltastar)$ is intensive by (C11); and $\Tbeta$ is
intensive because $\kappa$ and $\deltastar$ are deployment-class
constants independent of $|P|$.
\end{lemma}

\section{Proof}

\begin{proof}
We bound $h_{\mathrm{detect}}$ as the supremum of $\epsgap$ over
the detection window, apply (C11) per SPRT step, and bound
$T_{\mathrm{detect}}$ from above by the SPRT tail quantile.

\emph{Step 1 (per-step decomposition).}  Under the v3 boundary
convention, set $s_0 := t_0$ (the last safe SPRT step) and $s_n :=
t_0 + n$ in SPRT exposure clock.  For any
$0 \leq n \leq T_{\mathrm{detect}}$:
\[
  \epsgap(s_n) \;=\; \epsgap(s_0) + \sum_{i=1}^{n}
  (\Delta \epsgap)_i.
\]

\emph{Step 2 (apply (C11)).}  By Assumption~\ref{ass:C11}, each
per-step increment satisfies $(\Delta \epsgap)_n \leq
\rhogap(\deltastar)$.  Therefore:
\[
  \epsgap(s_n) \;\leq\; \epsgap(s_0) + n \cdot
  \rhogap(\deltastar).
\]

\emph{Step 3 (boundary condition).}  By the boundary convention,
$s_0 = t_0$ is the last SPRT step at which the deployment is still
within Layer~1's safe region.  Layer~1's intensivity result
(Lemma~\ref{lem:1}) gives directly:
\[
  \epsgap(s_0) \;=\; \epsgap(t_0) \;\leq\;
  h_{\mathrm{static}}(\theta).
\]
The first crossing event occurs at $s_1$, where (C11) bounds the
single-step jump by $\rhogap(\deltastar)$.  No buffer absorption is
required.

\emph{Step 4 (supremum over window).}  Combining Steps~1--3:
\[
  \sup_{0 \leq n \leq T_{\mathrm{detect}}} \epsgap(s_n)
  \;\leq\; h_{\mathrm{static}}(\theta) +
  T_{\mathrm{detect}} \cdot \rhogap(\deltastar).
\]
Therefore:
\[
  h_{\mathrm{detect}}(\theta) \;\leq\;
  h_{\mathrm{static}}(\theta) + \rhogap(\deltastar) \cdot
  T_{\mathrm{detect}}.
\]

\emph{Step 5 (apply Lemma 4 SPRT tail bound).}  By
Lemma~\ref{lem:4}, the SPRT detection time satisfies the
Wald--Hoeffding tail bound:
\[
  \Pr[T_{\mathrm{detect}} > T] \;\leq\; \exp(-\kappa T \deltastar).
\]
Setting $T = \Tbeta = (\log(1/\alpha) + \log(1/\beta)) /
(\kappa \deltastar)$ gives
$\Pr[T_{\mathrm{detect}} > \Tbeta] \leq \beta$, equivalently
$\Pr[T_{\mathrm{detect}} \leq \Tbeta] \geq 1-\beta$.

On the event $\{T_{\mathrm{detect}} \leq \Tbeta\}$ (probability
at least $1-\beta$):
\[
  h_{\mathrm{detect}}(\theta) \;\leq\;
  h_{\mathrm{static}}(\theta) + \rhogap(\deltastar) \cdot
  \Tbeta.
\]
This proves part~(ii) of the lemma.

\emph{Step 6 (lead-time-before-cascade composition).}  Separately
from the supremum bound, Lemma~\ref{lem:5d} composes Lemma~4 with
\citep[Phase~Redundancy]{lasser2026phase}'s metastable-lifetime
lower bound $T_{\mathrm{cascade}} \geq \tau_{\mathrm{meta}}$ to
yield $\Pr[\Tbeta \leq \tau_{\mathrm{meta}}] \geq 1 - \beta'$ with
$\beta' = \exp(-\kappa \tau_{\mathrm{meta}} \deltastar)$.  This
guarantees that detection completes before cascade with high
probability.  Importantly, $\tau_{\mathrm{meta}}$ enters as a
\emph{lower bound on cascade time} (correctly using
Phase~Redundancy's result), not as the integration horizon for
$\rhogap$.  This proves part~(iii).

\emph{Step 7 (intensivity).}  Each constant in the bound is
independent of $|P|$ over $\mathcal{D}$:
\begin{itemize}
\item $h_{\mathrm{static}}(\theta)$: intensive by Lemma~\ref{lem:1}.
\item $\rhogap(\deltastar)$: deployment-class constant by (C11).
\item $\Tbeta = (\log(1/\alpha) + \log(1/\beta)) /
  (\kappa \deltastar)$: intensive because $\alpha, \beta$ are
  monitor-design parameters, $\kappa$ is a sub-Gaussian constant of
  the SPRT log-likelihood (deployment-class via the channel-model
  structure), and $\deltastar$ is the union-class KL floor
  (Lemmas~\ref{lem:5b}, \ref{lem:5e}, deployment-class).
\end{itemize}
Therefore $h_{\mathrm{detect}}(\theta)$ is intensive in $|P|$ over
$\mathcal{D}$.  $\Box$
\end{proof}

\section{Discussion}

\paragraph{The two roles of cascade time.}
V3 distinguishes two roles previously conflated in v2:

\begin{itemize}
\item \emph{Integration horizon for $\rhogap$:} $\Tbeta$, the SPRT
  high-probability detection quantile, gives a legitimate
  \emph{upper} bound on $T_{\mathrm{detect}}$ derived from
  Lemma~4's Wald--Hoeffding tail.  This is what multiplies
  $\rhogap$ in the supremum bound.
\item \emph{Lead-time-before-cascade guarantee:}
  $\tau_{\mathrm{meta}} \leq T_{\mathrm{cascade}}$ (a \emph{lower}
  bound on cascade time, from
  \citep[Phase~Redundancy]{lasser2026phase}).  This is composed
  with $\Tbeta$ via Lemma~5d to give
  $\Pr[\Tbeta \leq \tau_{\mathrm{meta}}] \geq 1 - \beta'$, i.e.,
  detection completes before cascade.
\end{itemize}

V2 incorrectly used $T_{\mathrm{cascade}}$ for both roles, treating
a lower-bound quantity as if it were an upper bound on
$T_{\mathrm{detect}}$.  V3 separates them.

\paragraph{Lemma 5d integration item (Round B Q3 follow-up).}
At integration time, §4.10's Lemma~5d should be restated using the
union-class floor $\deltastar$ rather than $\deltaadv$ alone, to
match Layer~2's $\Aadv \cup \Aadv^{\mathrm{env}}$ class.  Round~B
recommended option~(c): define $\deltastar$ as the deployment-class
KL floor for the union, and restate Lemma~5d as the union-class
lead-time tail composition.  This is a small amendment to §4.10
and §A.4 that should accompany Lemma~6 integration.

\paragraph{Constant taxonomy.}
\begin{itemize}
\item $h_{\mathrm{static}}(\theta)$: intensive composition
  (Lemma~\ref{lem:1}).
\item $\rhogap(\deltastar)$: \emph{newly assumed Paper 10
  operational constant} via (C11); calibrated by per-channel
  sensitivity-to-gap measurement (Audit~7).
\item $\Tbeta = (\log(1/\alpha) + \log(1/\beta)) /
  (\kappa \deltastar)$: SPRT-derived detection quantile, intensive.
\item $\tau_{\mathrm{meta}}, T_{\mathrm{cascade}}$: source-derived
  cascade-time references from
  \citep[Phase~Redundancy]{lasser2026phase} (lower bound only;
  with the (C7) bounded co-evolution caveat).
\item $\beta$: monitor Type-II error parameter, design choice.
\item $\beta'$: lead-time-before-cascade tail, composed from
  Lemmas~5b/5e, 4, 5d.
\end{itemize}

\paragraph{Layer 2 operational form.}
Combining Lemma~6 with $\mathrm{Lip}(g) \leq K_{\mathrm{Lip}}$ from
(C6):
\[
  |g(\Ttruth) - g(\Pproxy)|(s) \;\leq\; K_{\mathrm{Lip}} \cdot
  \bigl(h_{\mathrm{static}}(\theta) + \rhogap(\deltastar) \cdot
  \Tbeta\bigr)
\]
with probability at least $1 - \beta$ (from Lemma~6 part~(ii)),
\emph{and} detection completes before cascade with probability at
least $1 - \beta'$ (Lemma~6 part~(iii)).  The two probabilistic
guarantees are composed via union bound: total failure probability
$\leq \beta + \beta'$.  The right-hand side is the sum of intensive
constants times $K_{\mathrm{Lip}}$, all $|P|$-independent.

\section{What this proof does and does not establish}

\textbf{Establishes:}
\begin{itemize}
\item Formal monitoring-clock formulation of (C11) with discrete
  (per-SPRT-step) and continuous variants.
\item (C11.CLK) clock-comparability sub-clause linking SPRT
  exposure clock to wall-clock cascade time.
\item Supremum-form $h_{\mathrm{detect}}$ definition matching
  Theorem 1's Layer 2 application.
\item Rigorous discrete-time boundary condition under the
  $t_0 := t_0^-$ convention.
\item Union-class KL floor $\deltastar$ for Layer 2's full
  detection class.
\item Intensive bound on $h_{\mathrm{detect}}$ as $h_{\mathrm{static}}
  + \rhogap \Tbeta$ with high probability $1-\beta$, where
  $\Tbeta$ is the SPRT detection quantile (legitimate upper
  bound).
\item Separate lead-time-before-cascade guarantee $\Tbeta \leq
  \tau_{\mathrm{meta}}$ with probability $1-\beta'$.
\end{itemize}

\textbf{Does not establish (deferred):}
\begin{itemize}
\item Tightness of $\rhogap(\deltastar)$.  Operators may calibrate
  tighter bounds for restricted threat models.
\item Existence of deployment configurations satisfying (C11) and
  (C11.CLK).  Verified empirically (Audit~7).
\item Numerical value of $\rhogap(\deltastar)$.  Calibrated
  deployment-by-deployment.
\item Structural derivation of (C11) from finer source-paper
  machinery.  Currently a Paper 10 operational assumption.
\end{itemize}

\section{Specific verification questions for v3 codex review}

\begin{enumerate}
\item Is the $T_{\mathrm{cascade}} \to \Tbeta$ replacement
  rigorous?  $\Tbeta = (\log(1/\alpha) + \log(1/\beta)) /
  (\kappa \deltastar)$ is a legitimate upper bound on
  $T_{\mathrm{detect}}$ via Lemma~4's Wald--Hoeffding tail.  Is
  the constant-tracking correct, and is the intensivity argument
  for $\Tbeta$ in Step~7 sound?  In particular, is $\kappa$ truly a
  deployment-class constant, or could it depend on $|P|$ via
  channel-model multiplicity?

\item Is the discrete boundary convention ($t_0 := t_0^-$ on the
  SPRT exposure clock) sufficient to make $\epsgap(t_0) \leq
  h_{\mathrm{static}}$ rigorous, or does it leak into Theorem~1's
  Layer~1$\to$Layer~2 transition statement (i.e., does the
  theorem-level proof need to adopt the same convention)?

\item Is (C11.CLK) the right framing for clock comparability, or
  should it be a top-level Theorem~1 condition (e.g., (C12)
  separate from (C11))?  The current draft treats (C11.CLK) as a
  sub-clause of (C11); is the conceptual coupling
  (gap-growth-rate-per-step + event-count-floor) tight enough to
  bundle them, or are they orthogonal deployment properties?

\item Is the union bound composition $\beta + \beta'$ in the
  Layer~2 operational form correct?  Or should the Layer~2 bound
  state the two events separately
  (Lemma~6(ii) and Lemma~6(iii)) and let Theorem~1 compose them
  with the appropriate union/intersection logic?

\item Is the constant taxonomy in the Discussion section now
  complete and accurate?  Are there any remaining hidden
  $|P|$-dependent constants in $\Tbeta$, $\rhogap$, or the SPRT
  parameters that we haven't surfaced?
\end{enumerate}

\end{document}
